% !TeX root = ../navier-stokes-manuscript.tex
% ============================================================
% 3.4 One Endpoint for the Whole Tower
% ============================================================
\subsection{One Endpoint for the Whole Tower}\label{sub:OneEndpointForTheWholeTower}

A smooth restart requires one value \(u(T_*)\) in every finite Sobolev space.
The complete spatial row gives the required space integrability, and the original equation converts it into enough time integrability to reach \(T_*\).

\begin{lemma}[The complete spatial row reaches every endpoint topology]\label{lem:CompleteSpatialRowEndpoint}
Let \((u,p)\) be the fixed pressure-normalized \NavierStokes{} history on \(I_*\).
If \(u\in\Szero(I_*)\), then, for every \(m\in\Nzero\),
\begin{equation}\label{eq:SectionThreeWOneOneEndpoint}
  u_t\in L^1(I_*;H^m),
  \qquad
  u\in W^{1,1}(I_*;H^m)
  \hookrightarrow C([0,T_*];H^m).
\end{equation}
\end{lemma}

\begin{proof}
The Leray projection rewrites the pressure-completed interaction as
\[
  (u\cdot\nabla)u+\nabla p
  =\Leray\nabla\cdot(u\otimes u).
\]
Under the fixed normalization, the Leray expression is exactly the combined transport-plus-pressure response.
Sobolev multiplication in three dimensions and the boundedness of the Leray and Riesz multipliers give
\[
  \norm{\Leray\nabla\cdot(u\otimes u)}_{H^m}
  \leq
  C_m\norm{u}_{H^{m+2}}^2.
\]
The complete spatial row and the finiteness of \(T_*\) yield
\begin{align*}
  \norm{\Leray\nabla\cdot(u\otimes u)}_{L^1(I_*;H^m)}
  &\leq
  C_m\norm{u}_{L^2(I_*;H^{m+2})}^{2}
  <\infty,\\
  \norm{\Delta u}_{L^1(I_*;H^m)}
  &\leq
  T_*^{1/2}\norm{u}_{L^2(I_*;H^{m+2})}
  <\infty.
\end{align*}
The momentum equation then places \(u_t\) in \(L^1(I_*;H^m)\).
More concretely, for \(0\leq s<t<T_*\),
\[
  \norm{u(t)-u(s)}_{H^m}
  \leq
  \nu(t-s)^{1/2}\norm{u}_{L^2((s,t);H^{m+2})}
  +C_m\norm{u}_{L^2((s,t);H^{m+2})}^{2}.
\]
Both terms tend to zero as \(s,t\uparrow T_*\), so the same velocity is Cauchy in \(H^m\) and reaches one endpoint there.
The inclusion \(L^2(I_*;H^m)\subset L^1(I_*;H^m)\) places \(u\) in the Bochner space \(W^{1,1}(I_*;H^m)\).
The Bochner fundamental theorem of calculus supplies the continuous representative in \eqref{eq:SectionThreeWOneOneEndpoint}.
\end{proof}

For each \(m\), let \(u_*^{(m)}\in H^m\) denote the endpoint value supplied by \eqref{eq:SectionThreeWOneOneEndpoint}.
If \(k>m\), the \(H^k\)-continuous representative, viewed in \(H^m\), and the \(H^m\)-continuous representative agree almost everywhere on \(I_*\).
Their continuity makes them identical on the closed interval.
Their terminal values are consequently compatible and define one element
\begin{equation}\label{eq:SectionThreeCommonEndpoint}
  u_*:=u(T_*)
  \in\bigcap_{m=0}^{\infty}H^m(\R^3)
  =\Hinfty(\R^3).
\end{equation}
No estimate uniform in \(m\) occurs in this conclusion.
Continuity of divergence from \(H^m\) to \(H^{m-1}\) makes \(u_*\) divergence-free.
The strong \(L^2\) trace gives \(u_*\) finite kinetic energy.

Apply the adjacent trace lemma to each temporal row in \eqref{eq:WholeTerminalAdjacentRows}.
It gives
\[
  \partial_t^N u\in C([0,T_*];H^m)
  \qquad(N,m\in\Nzero).
\]
For every \(N\in\Nzero\), compatibility across the spatial orders identifies these traces as one value
\begin{equation}\label{eq:SectionThreeVelocityJetTraces}
  V_N^*
  :=\lim_{t\uparrow T_*}\partial_t^N u(t)
  \in\Hinfty(\R^3).
\end{equation}
The differentiated equation now identifies these traces with the jet generated from \(u_*\).

\begin{proposition}[Pressure-complete endpoint jet]\label{prop:PressureCompleteEndpointJet}
The endpoint velocity \(u_*\) determines a unique sequence \((V_N^*,Q_N^*)_{N\geq0}\subset\Hinfty\times\Hinfty\) through
\begin{align}
  V_0^*&=u_*,\notag\\
  Q_N^*&=R_iR_j\sum_{a=0}^{N}\binom{N}{a}V_{a,i}^*V_{N-a,j}^*,
  \label{eq:SectionThreeEndpointPressureJet}\\
  V_{N+1}^*&=\nu\Delta V_N^*
  -\sum_{a=0}^{N}\binom{N}{a}(V_a^*\cdot\nabla)V_{N-a}^*
  -\nabla Q_N^*.
  \label{eq:SectionThreeEndpointVelocityJet}
\end{align}
This recursively generated sequence equals the left endpoint traces of the original velocity and pressure history.
\end{proposition}

\begin{proof}
The Fr\'echet space \(\Hinfty=\bigcap_mH^m\) is closed under spatial differentiation and multiplication.
Riesz transforms preserve every \(H^m\), so they preserve \(\Hinfty\).
Equations \eqref{eq:SectionThreeEndpointPressureJet} and \eqref{eq:SectionThreeEndpointVelocityJet} thus construct one unique pair at each finite order.

Fix a target Sobolev order \(m\).
The traces in \eqref{eq:SectionThreeVelocityJetTraces} converge in Sobolev spaces high enough for every product and derivative appearing in the recurrence.
Multiplication, spatial differentiation, and the Riesz transforms are continuous between those finite Sobolev spaces.
Passing \(t\uparrow T_*\) in \eqref{eq:SectionThreePressureRecurrence} and \eqref{eq:SectionThreeVelocityRecurrence} gives \eqref{eq:SectionThreeEndpointPressureJet} and \eqref{eq:SectionThreeEndpointVelocityJet}.
Induction in \(N\), beginning with \(V_0^*=u_*\), identifies every recursively generated coordinate with the corresponding trace of the history.
The arbitrariness of \(m\) gives the identities in \(\Hinfty\).
Spatial differentiation then identifies every mixed coordinate \(\partial_t^N\partial_x^\alpha u\) at the terminal face.
\end{proof}

The pressure limit has no new additive constant to choose at \(T_*\).
The decaying Riesz formula fixes \(Q_N(t)\) for every \(t<T_*\), and passage to the limit preserves that normalization.
Hence every left limit \(\partial_t^N\partial_x^\alpha u(T_*)\) is exactly the derivative obtained by inserting \(u_*\) into the pressure-normalized Navier--Stokes recurrences.
